\documentclass[11pt]{article}

\usepackage[margin=25mm]{geometry}

\usepackage{amsmath,amssymb,fontspec,unicode-math}

\IfFontExistsTF{DejaVu Serif}{\setmainfont{DejaVu Serif}}{\setmainfont{Latin Modern Roman}}

\IfFontExistsTF{Latin Modern Math}{\setmathfont{Latin Modern Math}}{\setmathfont{DejaVu Math TeX Gyre}}

\usepackage{xurl,hyperref}

\setlength{\parindent}{0pt}\setlength{\parskip}{6pt}

\begin{document}

\title{Prior-Guided Flow Matching for Target-Aware Molecule Design with Learnable Atom Number}

\author{ELVIA preliminary AI analysis}\date{21 October 2026}\maketitle

Preliminary assessment: Minor.\par

Independent scientific review remains pending. Author adoption has not been requested.

The released categorical update uses a mode rather than the declared categorical draw. The affinity argument also mistakes positive Gaussian-noise improvement for proximity to a local minimum.

\section*{Review scope}

Rapid current-version reassessment of every substantive concern in the supplied LaTeX report: the implicated paper passages and relevant author-released source paths were inspected, existing proposed corrections were briefly checked, and impact was calibrated to the demonstrated scope. No author-provided or historical verification program was executed.

\section*{Problem 1: The released discrete state update selects a mode instead of sampling the declared categorical path}

Source: Section 4.1, Eqs. (6)–(8) and Algorithm 2 final discrete step; uniform initialization, PDF p. 5.

Assessment: supported; impact: minor

\subsection*{Source claim (paraphrase)}

Paraphrase of the manuscript: Under the conditional path in Eq. (6), an atom type has categorical probabilities \(q_\alpha\). The vector-field construction in Eq. (7), its update in Eq. (8), and Algorithm 2’s final discrete step provide the sampling procedure for this path. Initialization uses uniformly distributed atom types (Section 4.1, page 5). Algorithm 2, line 13, page 23, explicitly calls for sampling from the updated categorical distribution.

The relevant public implementation is models/molopt\_score\_model\_guide.py, function sample\_guided\_flow\_VP, beginning at line 659. Lines 673–674 initialize the accumulator with the one-hot initial category. Lines 726–729 sample a predicted clean endpoint using torch.multinomial. Lines 732–735 update the accumulator. Lines 737–738 apply a softmax to that accumulator and then call argmax to choose the next atom type. The wrapper scripts/sample\_flow\_VP\_guide.py initializes categories uniformly; the atom-type guidance weight defaults to zero. These locations concern the guided sampler actually being examined, rather than a similarly named function in the unguided model file.

\subsection*{Analysis}

Let the initial category \(a_0\) be uniform on \(\{e_1,e_2\}\), the endpoint be \(e_1\), and \(\alpha=1/4\). The conditional probabilities are \((1-\alpha)a_0+\alpha e_1\). Selecting their mode gives \(e_1\) with marginal probability \(1/2\), while drawing categorically gives \((1+\alpha)/2=5/8\). Earlier endpoint sampling cannot turn the later argmax into the declared categorical draw.

\subsection*{Finding and scope}

The current guided source still uses the mode. The exact conditional-oracle reduction disproves equivalence to the declared intermediate categorical law, without reproducing a full trained trajectory.

\subsection*{Proposed correction}

For the fixed-endpoint construction, draw from \((1-\alpha)a_0+\alpha a_1\); averaging over uniform \(a_0\) gives exactly the declared conditional marginal. A complete learned, guided sampler still needs a derivation ensuring valid probabilities; sampling softmax of an already normalized probability vector is not equivalent.

A local correction or qualification is proposed where specified; author adoption is not established.

\section*{Problem 2: Gaussian-noise improvement does not imply proximity to an affinity minimum}

Source: Appendix A.4, Eqs. (30)–(33), arXiv v2 PDF p. 22.

Assessment: supported; impact: minor

\subsection*{Source claim (paraphrase)}

Appendix A.4, PDF page 22, Eqs. (30)–(33), compares deterministic affinity \(f(n)\) with \(E[f(n+\delta Z)]\). A second-order expansion yields approximately

\[ E[f(n+\delta Z)]\approx f(n)+\frac{\delta^2}{2}f''(n). \]

The following prose infers that the predicted count lies near a local affinity minimum from the observed benefit of Gaussian noise.

\subsection*{Analysis}

Let \(\Phi\) be the standard normal CDF and \(f(x)=\Phi(x-1)\), a bounded smooth strictly increasing function with no local minimum. For \(n=1/2\) and \(Z\sim N(0,1)\), Gaussian convolution gives \(\mathbb E f(n+\delta Z)=\Phi(-1/(2\sqrt{1+\delta^2}))>\Phi(-1/2)=f(n)\) whenever \(\delta>0\). Noise improvement can therefore arise from curvature without proximity to a minimum.

\subsection*{Finding and scope}

The bounded monotone counterexample removes the minimum interpretation while allowing the reported empirical noise benefit to remain true.

\subsection*{Proposed correction}

Retain a conditional curvature statement: if \(f\in C^4(\mathbb R)\) with \(|f^{(4)}|\le M\), then \(\mathbb E f(n+\delta Z)-f(n)=\delta^2 f^{\prime\prime}(n)/2+r\), with \(|r|\le M\delta^4/8\). This does not imply stationarity or a nearby minimum, nor automatically describe the clipped and rounded integer-count pipeline.

A local correction or qualification is proposed where specified; author adoption is not established.

\section*{Limits}

No author-provided program, training procedure, model inference or wet-lab experiment was executed in this review.\par

Historical execution results are retained only with that attribution; a source-level defect alone does not identify the program or configuration used for the published experiments.\par

No author contact, adoption, independent human scientific review or novelty determination is established.\par

The findings apply only to the identified claims, source versions and stated assumptions; they do not certify or invalidate the entire paper.\par

\section*{Source identity}

arXiv:2509.01486v2 (4 November 2025)

\url{https://arxiv.org/pdf/2509.01486v2}

Reviewed PDF SHA-256: \nolinkurl{a5c4992b4c3ab258c0f000e6a2f6941a436d9eb7b710cfe61dd34e9230eea7c6}

The exact inspected current paper PDF is identified by its SHA-256. The current arXiv paper matches the four authors and retains both implicated arguments; its appendix pagination differs from the proceedings. Current PAFlow head remains the pinned release. An exact PDF acquisition timestamp is not established in this reassessment; source-currentness checks have their own recorded timestamps and file modification times were not substituted.

\end{document}
